\chapter*{Discussion}
\addcontentsline{toc}{chapter}{Discussion}
\label{chap:discussion}

\section*{Introduction}

Cette discussion interprète les résultats du pilote historique, les compare aux travaux du chapitre~\ref{chap:etat-art}, examine ce qui menace la validité de l'étude et en tire quelques enseignements pour la conception de contrôleurs décentralisés. Le périmètre est clair. Nous parlons du pilote, pas de l'implémentation actuelle, dont les entraînements et les campagnes comparatives restent à faire. Les implications pratiques et la suite expérimentale terminent la discussion, avant la conclusion générale.

\section*{Interprétation des résultats}

Ce que le pilote établit d'abord, c'est que la chaîne historique fonctionne. Le simulateur, le contrôleur à messages, la fonction de perte et l'optimisation produisent ensemble une configuration circulaire plus proche de la cible que l'état initial. Cela ne dit rien d'une convergence asymptotique, qui reste hors de portée. Par construction, les paramètres partagés appliquent le même réseau à chaque robot, et l'agrégation locale absorbe un nombre variable de voisins. Le protocole indique aussi une meilleure adéquation au cercle que la règle de regroupement implémentée.

Il ne faut pourtant pas faire dire plus à cette observation. La baseline de regroupement n'utilise pas la même consigne que le GNN, et elle n'est ni apprise ni ajustée à la tâche de cercle. L'écart de distance de Chamfer ne permet donc pas de répondre à la question de savoir si le GNN est meilleur à information égale. La question réellement posée est plus modeste, et le pilote y répond. Dans la configuration retenue, la politique entraînée se rapproche davantage du cercle qu'une règle qui cherche seulement à regrouper les robots.

Un point de protocole mérite d'être précisé, car il change la lecture. Le pilote évalue le modèle issu du dernier niveau du curriculum réduit, pas un modèle entraîné sur le seul cercle. Ce choix visait à évaluer le contrôleur réellement destiné à la démonstration, mais il ajoute une variable. Les résultats reflètent l'effet conjoint du curriculum et de l'architecture. Une campagne définitive devra les séparer, en évaluant à la fois le modèle spécialisé sur le cercle et celui du curriculum, sur les mêmes épisodes.

Après retrait de 20\,\% des robots, la distance moyenne se dégrade peu. C'est un signal à étudier, pas une preuve de robustesse générale. Le taux de succès, nul à ce niveau de pannes avec le seuil retenu, rappelle que la moyenne de Chamfer ne résume pas toute la performance. À vingt robots, le premier franchissement admissible se produit rarement. Or le champ historique appelé « pas de convergence » n'exige pas de rester sous le seuil ensuite, et il ne dit pas, à lui seul, si la formation s'est stabilisée. La mesure actuelle, fondée sur un suffixe sous le seuil, obéit à une autre convention et reste limitée à un horizon fini.

\section*{Comparaison avec les travaux antérieurs}

Il est utile de situer ces résultats par rapport aux travaux du chapitre~\ref{chap:etat-art}, en séparant ce qui tient à l'architecture de ce qui tient au protocole.

Tolstaya et al. montrent qu'un contrôleur GNN peut être entraîné pour des tâches de robotique en essaim sur un graphe de communication mobile \cite{tolstaya2020}. Notre travail appartient à la même famille d'architectures, mais la tâche change. Il faut converger vers une géométrie de cible donnée, et non suivre une trajectoire ou éviter des obstacles. Un objectif de forme impose de couvrir le contour, ce qui n'a pas d'équivalent direct en navigation.

Gama et al. étudient les GNN comme contrôleurs décentralisés et explorent le transfert vers des équipes plus grandes ou des scénarios nouveaux, notamment par imitation \cite{gama2020,gama2022}. Le passage à l'échelle observé dans le pilote va dans le même sens que leurs conclusions, puisque l'architecture ne s'effondre pas quand le nombre de robots change. La différence tient à l'apprentissage. Ici, le signal vient d'une perte de forme évaluée en fin d'épisode et remontée par simulation différentiable, et non de l'imitation d'un expert. On évite ainsi de dépendre d'un contrôleur de référence, en payant un horizon de rétropropagation égal à la durée de l'épisode.

Jiang, Huang et Guo proposent une formation décentralisée de bout en bout fondée sur un GNN \cite{jiang2023}. Le rapprochement est direct, mais les hypothèses de perception diffèrent. Leurs observations viennent de capteurs embarqués, transformées en cartes locales, alors que notre prototype travaille sur des états simulés avec une consigne géométrique donnée avant l'épisode. Comparer les performances n'a de sens qu'à hypothèses de perception comparables.

Les contrôleurs écrits à la main pour générer des formes \cite{hsieh2008,li2019} restent la référence à atteindre. Ils sont interprétables et ne demandent aucun entraînement, mais il faut les concevoir pour chaque famille de formes. Nous ne prétendons pas les remplacer. Nous montrons seulement qu'une loi apprise, exécutée localement, peut être entraînée de façon reproductible sur un ensemble de formes hétérogènes. La comparaison équitable, avec une règle manuelle qui recevrait une consigne de contour comparable, n'a pas encore été faite, et c'est la principale dette expérimentale du travail.

\section*{Menaces à la validité}

\subsection*{Validité interne}

Le premier facteur limitant est le nombre de répétitions. Avec dix épisodes par condition, on peut valider une chaîne de mesure et repérer de gros écarts, mais pas estimer précisément une moyenne, encore moins comparer des distributions. Un facteur secondaire, bien réel, est le couplage entre architecture et protocole. Le modèle évalué vient d'un curriculum, ce qui brouille l'attribution des résultats au seul passage de messages.

Le deuxième facteur est l'absence, dans les résultats du pilote, d'une baseline apprise sans communication. Cette variante existe maintenant dans l'implémentation, mais elle n'a pas encore été entraînée ni évaluée dans une campagne actuelle, et l'écart historique entre le GNN et le regroupement reste donc un écart entre un système conditionné par une cible et un système qui l'ignore. Enfin, avec une seule graine d'entraînement, les résultats décrivent le comportement d'un entraînement particulier, pas la performance attendue de l'architecture.

\subsection*{Validité externe}

Les robots du simulateur partagent un repère commun, connaissent leurs positions exactes et communiquent sans délai ni perte. Les collisions sont pénalisées par une fonction continue, sans dynamique de contact. Un robot en panne disparaît proprement du graphe, alors qu'un défaillant réel peut rester immobile, faire obstacle ou transmettre des valeurs aberrantes. L'arène est finie et ses bords modifient les trajectoires, si bien qu'aucun résultat ne s'extrapole à un espace ouvert. Les épisodes sont aussi très courts, vingt pas dans le pilote, bien en deçà d'une mission réelle.

Ces hypothèses rendent l'étude contrôlable, mais elles en limitent la portée. Les résultats établissent une faisabilité en simulation 2D, pas un transfert vers des plateformes physiques.

\subsection*{Validité de construit}

Les mesures elles-mêmes comportent des approximations. La distance de Chamfer tient compte des deux ensembles mais ignore la lisibilité d'un symbole, et deux configurations peuvent avoir des distances voisines alors que l'une reste interprétable et l'autre non. Le seuil de succès, fixé à 0,5 dans le pilote, est un choix opérationnel dont dépend directement le pourcentage de réussite. Le taux de collisions mesure la fréquence de rapprochements sous une distance de sécurité, pas des impacts mécaniques. Quant à la régularité des voisins, elle détecte des amas sans garantir une couverture complète du contour.

Cela ne rend pas ces mesures inutiles. Cela oblige à les présenter ensemble, avec leur définition, plutôt que de retenir celle qui donne le résultat le plus flatteur.

\section*{Implications pour la conception de contrôleurs décentralisés}

Trois enseignements se dégagent du travail, sous réserve des limites ci-dessus.

D'abord, le prototype qui combine consigne et messages a produit des formations non triviales dans les scénarios étudiés. On ignore encore laquelle des deux contributions pèse le plus, et les données disponibles ne les séparent pas. C'est justement l'objet de l'ablation sans communication et de la comparaison entre identifiant seul et géométrie, toutes deux préparées.

Ensuite, la moyenne évite que la somme des messages croisse mécaniquement avec le nombre de voisins. Elle ne garantit pas pour autant que leurs distributions restent identiques quand la densité change. Les résultats historiques à 40 et 100 robots sont compatibles avec ce choix de conception, sans en établir l'effet causal. Une attention à poids normalisés peut aussi former une moyenne pondérée, et ses propriétés comme sa performance méritent d'être évaluées plutôt que rejetées d'emblée.

Enfin, la progression cumulative est censée limiter l'oubli. Le pilote n'isole pas son gain par rapport à un entraînement non progressif, et conserver les anciennes cibles ne dispense pas de mesurer leur performance individuelle à chaque niveau, sur des épisodes appariés. Le tableau des pertes par cible du chapitre~\ref{chap:experiences} est un premier pas dans cette direction.

\section*{Considérations pratiques et sociales}

On justifie souvent les essaims robotiques par leur faible coût unitaire et leur redondance, et des plateformes comme les Kilobot montrent qu'on peut déployer des centaines de robots simples \cite{rubenstein2014}. Cela change aussi les enjeux d'évaluation. Un contrôleur qui marche en simulation mais repose sur des hypothèses fragiles de repère ou de communication peut imposer, à l'échelle du groupe, des coûts de recalibrage considérables.

Un second point touche à l'usage de la simulation. Entraîner et valider numériquement avant tout déploiement physique réduit le nombre d'essais sur matériel, donc la consommation et les pertes qui vont avec. La reproductibilité des campagnes, obtenue ici par des configurations et des graines explicites, y contribue, car rejouer une expérience coûte du temps de calcul alors que régénérer des robots coûte des matières.

Enfin, une formation de forme par essaim est un dispositif visible dans l'espace public. Toute extension vers des plateformes réelles devra traiter le bruit, l'occupation de l'espace et le consentement des personnes concernées, ainsi que les usages ambivalents d'un tel dispositif. Nous n'abordons pas ces questions expérimentalement dans ce mémoire, et elles figurent parmi les limites de périmètre.

\section*{Suite expérimentale}

Les campagnes préparées prévoient au moins 50 départs par condition et cinq graines d'entraînement, avec une séparation nette des deux sources de variabilité. Elles incluent le bruit d'observation, la suppression aléatoire d'arêtes et la réduction de portée. Supprimer des arêtes n'équivaut pas à bruiter le contenu des messages. La comparaison prioritaire oppose la politique sans communication à son équivalent avec messages, à consigne identique. Les modules distinguent désormais proximité, couverture, régularité et sécurité, et un questionnaire en aveugle prépare l'étude de lisibilité. Leur existence ne remplace ni l'exécution des campagnes ni la collecte de réponses humaines. Le Tableau~\ref{tab:controller-matrix} récapitule, contrôleur par contrôleur, l'information disponible et le statut de chaque comparaison.

\begin{table}[htbp]
  \centering
  \caption[Information reçue par les contrôleurs actuels]{Information reçue par les contrôleurs actuels (voisinage, identifiant et géométrie). Les variantes neuronales doivent encore être entraînées et comparées, et cette matrice ne représente pas des performances.}
  \label{tab:controller-matrix}
  \begin{tabular}{lcccl}
    \toprule
    \textbf{Contrôleur} & \textbf{Voisins} & \textbf{ID forme} & \textbf{Géométrie cible} & \textbf{Poids ajustables} \\
    \midrule
    Regroupement (flocking) & Oui & Non & Aucune & Non \\
    Règle manuelle & Oui & Non & Déplacement au plus proche (direction + distance) & Non \\
    GNN identifiant seul & Oui & Oui & Aucune & Oui \\
    GNN avec géométrie & Oui & Oui & Direction unitaire vers le plus proche & Oui \\
    GNN sans communication & Non & Oui & Direction unitaire vers le plus proche & Oui \\
    \bottomrule
  \end{tabular}
  
  \vspace{0.5em}
  \footnotesize
  Les poids ajustables sont les paramètres à apprendre, sans entraînement actuel attesté.\\
  \textbf{Le pilote n'a évalué que l'ancien GNN et le regroupement.} Les autres variantes sont nouvellement préparées, sans entraînement actuel. Le GNN sans communication conserve ses 2 tours, mais son graphe et ses agrégats sont nuls.
\end{table}

Un premier axe de raffinement concerne l'agrégation. Comparer la moyenne à une agrégation tronquée ou pondérée permettrait de mesurer la sensibilité du contrôleur à un voisin défaillant. Un second axe concerne la profondeur. Comparer un, deux et trois tours à budget d'entraînement égal évalue l'effet conjoint de la portée informationnelle, de la profondeur et de la capacité paramétrique, et les séparer demandera des contrôles supplémentaires.

À plus long terme, il faudra remplacer le repère commun par des observations relatives, une localisation distribuée ou une balise mesurable. Passer d'une dynamique à vitesse directe à un modèle cinématique ou dynamique plus réaliste permettra ensuite d'évaluer le transfert vers des robots physiques.

\section*{Conclusion partielle}

Cette discussion a replacé les observations du pilote historique dans la littérature et précisé leurs limites, qui tiennent à l'asymétrie de consigne, au peu de répétitions, à la dépendance à une graine d'entraînement et à une simulation simplifiée. Elle a séparé l'exécution à taille variable de la généralisation mesurée, et la proximité géométrique de la lisibilité. La conclusion générale reprend les apports de l'étude et les validations expérimentales qui restent à conduire.

\chapter*{Conclusion générale et perspectives}
\addcontentsline{toc}{chapter}{Conclusion générale et perspectives}
\label{chap:conclusion}

Ce travail a mis en place une chaîne expérimentale reproductible pour étudier la formation de formes par un essaim piloté par GNN. Le système réunit un simulateur 2D, des formes représentées par des nuages de points, une politique locale à paramètres partagés et un apprentissage différentiable. Le cercle a offert une première observation, visuelle et chiffrée, du fonctionnement du prototype historique, mais cette observation ne valide pas les performances de l'implémentation actuelle. Génération de données, entraînement, évaluation et visualisation sont des modules séparés, de sorte que les expériences se rejouent à partir de leurs configurations.

Le résultat principal n'est pas encore une démonstration sur robots physiques. Il est plus précis. Dans le simulateur retenu, une politique GNN partagée peut recevoir des informations locales, échanger des messages avec ses voisins et produire des mouvements qui rapprochent l'essaim d'une cible. Le modèle peut aussi s'exécuter avec un nombre de robots différent de celui de l'entraînement. Cela découle du partage des paramètres et de l'agrégation des messages, mais la qualité obtenue doit être vérifiée expérimentalement pour chaque scénario.

Le protocole pilote suggère par ailleurs que le contrôleur appris s'ajuste mieux à une cible circulaire qu'une règle locale de regroupement, y compris dans les scénarios de retrait étudiés. Nous restreignons volontairement cette conclusion, puisque nous n'avons que dix graines par condition et une baseline qui ne reçoit pas la même consigne. Pour les formes symboliques, la couverture n'est encore qu'approximative. La chaîne de traitement accepte plusieurs types de contours, mais rien ne permet d'affirmer qu'elle forme des lettres ou des nombres lisibles et généralisables.

Les premières perspectives sont expérimentales. Il faut augmenter le nombre de répétitions, fixer un jeu de formes inconnues, mesurer des intervalles de confiance et garder les épisodes défavorables autant que les favorables. Il faut aussi entraîner et évaluer les variantes déjà implémentées, politique sans communication et plusieurs profondeurs de passage de messages, puis les confronter à la règle manuelle. La consigne de cette règle reste plus riche, puisqu'elle conserve la longueur du déplacement vers la cible, et l'équivalence d'information doit donc être explicitée et non supposée. Ces comparaisons permettront d'interroger l'apport propre des messages.

Les perspectives techniques touchent au réalisme du simulateur. Des bruits de position et de communication, des délais, des obstacles et une dynamique de robot non holonome rapprocheraient le cadre de situations réelles. À plus long terme, la consigne géométrique pourrait venir d'une perception locale, d'une balise ou d'une carte distribuée. Passer à une plateforme physique demanderait alors une estimation de pose, une stratégie de sécurité et une validation progressive, d'abord sur quelques robots, puis sur un essaim plus grand.

Le projet répond donc en partie à la question initiale. Il montre qu'un contrôleur GNN local peut former des formes en simulation, et il fournit les outils d'une évaluation plus rigoureuse. La suite devra établir, avec des expériences mieux contrôlées, dans quelles conditions cette approche apporte un avantage réel sur les méthodes classiques de contrôle distribué.